\begin{center}
\Huge
Stamfunktioner og ubestemte integraler
\end{center}

\section*{Stamfunktioner}
\stepcounter{section}

Hvis vi differentierer en differentiabel funktion $f$, så kan vi spørge os selv, om vi kan gå den anden vej - givet et funktion $f'$, kan vi så bestemme $f$? Dette kan vi godt op til konstantled. 
\begin{defn}[Stamfunktion]
En differentiabel funktion $F$ kaldes en stamfunktion til en funktion $f$, hvis det gælder, at 
\begin{align*}
F'(x) = f(x).
\end{align*}
\end{defn}
\begin{exa}
	Vi betragter funktionerne $F$ og $f$ givet ved henholdsvis
	\begin{align*}
		F(x) = x^2 + 5
	\end{align*}
	og 
	\begin{align*}
		f(x) = 2x.
	\end{align*}
	Vi undersøger, om $F$ er en stamfunktion til $f$ ved at differentiere:
	\begin{align*}
		F'(x) = 2x + 0 = 2x = f(x),
	\end{align*}
	så $F$ er derfor en stamfunktion til $f$, da $F' = f$. 
	Tilsvarende må der gælde, at funktionen $G$ givet ved
	\begin{align*}
		G(x) = x^2 + 10
	\end{align*}
	også er en stamfunktion til $f$, da
	\begin{align*}
		G'(x) = 2x + 0 = 2x = f(x).
	\end{align*}
\end{exa}
Af eksemplet kan vi se, at der ikke er en entydig stamfunktion til en funktion. Når vi differentierer, går konstanter ud, så vi kan allerhøjst finde stamfunktioner op til en konstant forskel. Dette leder os til følgende definition for det
ubestemte integral.
\newpage
\begin{defn}[Ubestemt integral]
	Det ubestemte integral for en kontinuert funktion $f$ defineres som funktionen
	\begin{align*}
		\int f(x) dx = F(x) + k,
	\end{align*}
	hvor $F$ er en stamfunktion til $f$ og $k\in \mathbb{R}$ er en konstant.
\end{defn}
\begin{exa}
	Vi får funktionen $f(x) = x^4$ givet og ønsker at bestemme det ubestemte integral
	\begin{align*}
		\int f(x) dx = \int x^4 dx = \frac{1}{5}x^5 + k.
	\end{align*}
	Vi kan tjekke efter og får
	\begin{align*}
		\left(\int f(x) dx \right)' = \left(\frac{1}{5}x^5+ k\right)' = \frac{1}{5}\cdot 5x^4 = x^4, 
	\end{align*}
	så $\frac{1}{5}x^5 + k$ er faktisk en stamfunktion til $x^4$. 
\end{exa}
Vi har følgende sætning, der fortæller 
\begin{setn}
Lad $F_1$ og $F_2$ være stamfunktioner til en funktion $f$. Så afviger $F_1$ og $F_2$ højst med en konstant. Mere præcist gælder der, at 
\begin{align*}
F_1-F_2 = k.
\end{align*}
\end{setn}
\begin{proof}
Vi antager, at $F_1$ og $F_2$ begge er stamfunktioner til $f$, og vi differentierer så $F_1-F_2$ og får
\begin{align*}
	(F_1(x)-F_2(x))' = F'_1(x)-F_2'(x) = f(x)-f(x) = 0.
\end{align*}
Da $(F_1-F_2)' = 0$, så må $F_1-F_2=k$, da kun konstanter differentieres til $0$.  
\end{proof}

Vi kan som det var tilfældet med afledede funktioner lave en tabel med stamfunktioner. Denne kan ses på Tabel \ref{tab:stamfunktioner}.
\begin{table}[H]
	\centering
	\rowcolors{1}{Teal!40}{Teal!20}
	\begin{tabular}{cc}
		$f(x)$ & $F(x)$ \\
		$k$    & $kx$ \\
		$1$    & $x$ \\
		$x$    & $\frac{1}{2}x^2$ \\
		$x^a$  & $\frac{1}{a+1}x^{a+1}$ \\
		$\frac{1}{x}$ & $\ln(|x|)$ \\
		$\cos(x)$ & $\sin(x)$ \\
		$\sin(x)$ & $-\cos(x)$ \\
	\end{tabular}
	\caption{Stamfunktioner for $f$, hvor konstantleddet er udeladt.}
	\label{tab:stamfunktioner}
\end{table}

\newpage

\section*{Stamfunktioner gennem et punkt}
Vi kan godt forestille os en situation, hvor vi får en betingelse på stamfunktionen $F$, så vi kan bestemme en bestemt 
værdi for $k$. Vi betragter følgende eksempel.

\begin{exa}
	Vi ønsker at bestemme en stamfunktion til funktionen $f$ givet ved
	\begin{align*}
		f(x) = \frac{1}{x}, x > 0,
	\end{align*}
	der går gennem punktet $(1,7)$. Vi bestemmer først det ubestemte integral.
	\begin{align*}
		F(x) = \int \frac{1}{x} dx = \ln(x) + k
	\end{align*}
	Vi indsætter så punktet i stamfunktionen og får
	\begin{align*}
		F(1) = \ln(1) + k = 0 + k.
	\end{align*}
	Vi ved, at denne skal være lig $7$ i punktet, så vi må have, at $k = 7$ og stamfunktionen er derfor givet ved
	\begin{align*}
		F(x) = \ln(x) + 7.
	\end{align*}
\end{exa}



\subsection*{Opgave 1}
Afgør, om følgende funktioner $F$ er stamfunktioner til funktionerne $f$.
\begin{enumerate}[label=\roman*)]
	\item $F(x) = x^2 + 10x + 10$ og $f(x) = 2x+10$.
	\item $F(x) = x^4$ og $f(x) = x^3$.
	\item $F(x) = 2\ln(x)$ og $f(x) = \frac{2}{x}+4$.
	\item $F(x) = x^5+x^4+x^3+x^2+x+1$ og $f(x) = \frac{1}{5}x^4+\frac{1}{4}x^3+\frac{1}{3}x^2+1$.
	\item $F(x) = 10x$ og $f(x) = 10$.
	\item $F(x) = 5$ og $f(x) = 0$.
	\item $F(x) = e^{x^2}$ og $f(x) = 2xe^{x^2}$.
	\item $F(x) = \ln(x)e^{2x}$ og $f(x) = e^{2x}(\frac{1}{x}+2\ln(x))$.
\end{enumerate}
\subsection*{Opgave 2}
Bestem følgende ubestemte integraler
\begin{align*}
	&1) \  \int 10 dx   &&2) \  \int x dx  \\
	&3) \  \int -12x dx   &&4) \ \int 2x^2 dx    \\
	&5) \  \int 3x^2+2x+1 dx &&6) \ \int \frac{10}{x} dx     \\
	&7) \  \int 5^x+10x dx &&8) \ \int e^{3x} dx     \\
	&9) \   \int 2e^{2x} dx  &&10) \  \int 2xe^{x^2} dx    \\
\end{align*}
\subsection*{Opgave 3}
Bestem en stamfunktion til følgende funktioner $f$, der går gennem det givne punkt $P$. 
\begin{enumerate}[label=\roman*)]
	\item $f(x) = 3x^2$, $P(2,5)$.
	\item $f(x) = x^2 + 2x - 1$, $P(-1, 10)$.
	\item $f(x) = \cos(x)$, $P(0,4)$.
	\item $f(x) = \frac{1}{x}+x^2$, $P(1,3)$.
	\item $f(x) = \frac{1}{2\sqrt{x}}$, $P(4,-2)$.
\end{enumerate}

\subsection*{Opgave 4}
Bestem en forskrift for $f$, hvis 
\begin{enumerate}[label=\roman*)]
	\item $f''(x) = 12x^2 - 6x$, $f'(1) = 2$, $f(2) = 1$.
	\item $f''(x) = \cos(x)$, $f'(0) = 5$, $f(0) = 3$.
	\item $f''(x) = x^{-1}$, $f'(2) = 1$, $f(2) = 2$.
	\item $f''(x) = 2x^2-x+4$, $f'(4) = 10$, $f(-2) = 1$.
\end{enumerate}